\label{exp:cem03}

\subsubsection*{Постановка}
\textbf{Физический вопрос.}
Стандартная модель фиксирует массы фермионов, бозона Хиггса и масштаб $v$ как измеренные величины.
Модель утверждает, что те же числа выводятся из лестницы масштабов и bare-параметров (\cref{thm:vev,prop:ladder-scale-steps}) без подгонки каждой массы отдельно.

\textbf{Что проверяем.}
\begin{enumerate}[label=\textbf{(\alph*)},leftmargin=*,itemsep=0.25em]
\item $v$ и массы $m_H$, $m_p$, $m_e$, $m_n$ из формул главы SI согласуются с PDG в полосе модели.
\item Bare $\sin^2\theta_W^{(\mathrm{bare})}=3/13$ и $\alpha_s(M_Z)$ из seed (\cref{prop:alpha-s-seed}) остаются в согласованном блоке с массами.
\item Сводная таблица модель/PDG (\cref{tab:si-pdg-ce-m}, глава~\cref{ch:si-sm}) не смешивает зависимые якоря как независимые тесты (\cref{sec:multiple-observables}).
\end{enumerate}

\textbf{Рабочая гипотеза.}
Относительные отклонения масс остаются $\mathcal{O}(10^{-3})$--$\mathcal{O}(10^{-4})$ при фиксированной лестнице; ни одна строка сводной таблицы не выходит за verify-полосу без объяснения upstream.

\textbf{Наблюдаемые.}
Строки сводной таблицы PDG; контрольные $|\delta|$ по отдельным массам (табл.~\ref{tab:exp-ce-m-03-checks}).

\textbf{Критерий согласия с теорией.}
Каждая масса в полосе, зафиксированной для SM-suite; интерпретация только с учётом графа зависимостей якорей.

\subsubsection*{Теоретическая часть}
Теорема о $v$ и шаги лестницы задают дискретные масштабы; массы частиц --- функции этих масштабов и bare coupling (\cref{ch:si-sm,ch:alpha}).
Сравнение с PDG --- проверка \emph{всего блока} предсказаний, а не «одна масса = один степень свободы''.

\subsubsection*{Ход эксперимента}
Воспроизведение: \texttt{python verify\_principles.py --suite sm}.
Сводка модель/PDG: табл.~\ref{tab:si-pdg-ce-m} (глава~\cref{ch:si-sm}, приложение).

\begin{table}[htbp]
\centering
\small
\caption{Относительные отклонения масс к PDG (CE-M-03)}
\label{tab:exp-ce-m-03-checks}
\begin{tabular}{@{}lll@{}}
\toprule
Наблюдаемая & Согласие с полосой & $|\delta|$ \\
\midrule
$m_H$ & да & $4{,}73\times 10^{-4}$ \\
$m_p$ & да & $3{,}16\times 10^{-3}$ \\
$m_e$ & да & $4{,}50\times 10^{-3}$ \\
$m_n$ & да & $3{,}44\times 10^{-3}$ \\
\bottomrule
\end{tabular}
\end{table}

\subsubsection*{Анализ, расчёт погрешности и выводы}
Все четыре массы в одном порядке отклонения от PDG; это согласовано с одной лестницей, а не с четырьмя независимыми fits.
Строки табл.~\ref{tab:si-pdg-ce-m} читаются вместе с \cref{sec:multiple-observables}: повторное использование $v$ или $m_e$ в разных строках не удваивает статистический вес.

Погрешность: тип~B (сравнение с PDG как внешним эталоном).

\textbf{Вывод.}
Лестница масс и bare-параметры дают SM-спектр в заявленной полосе без построчной подгонки.
\textbf{Статус записи:} опыт закрыт по сводной таблице и табл.~\ref{tab:exp-ce-m-03-checks}.
